\section*{Capítulo \ref{ch:g8:voltage-voltmeter} --- La tensión y el voltímetro}

\begin{solution}{exo:g8:voltage-voltmeter:1}
El empujón eléctrico disponible entre dos puntos --- como la altura de
una cascada entre su cima y su base moviendo la rueda del molino. Al
ser una diferencia, necesita dos puntos; su \omterm{def:g6:measurement-in-science:measuring}{unidad} es el
\omterm{def:g8:voltage-voltmeter:voltage}{voltio}, \unit{V}.
\end{solution}

\begin{solution}{exo:g8:voltage-voltmeter:2}
\qty{1.5}{V}; \qty{4.5}{V}; \qty{9}{V}; \qty{12}{V}; \qty{230}{V}.
\end{solution}

\begin{solution}{exo:g8:voltage-voltmeter:3}
Ninguna corriente --- no hay lazo. El \omterm{def:g8:voltage-voltmeter:voltmeter}{voltímetro} lee los
\qty{1.5}{V} enteros: \omterm{def:g8:voltage-voltmeter:voltage}{tensión} sin corriente --- el empujón puede
esperar; el desfile no puede empezar sin él.
\end{solution}

\begin{solution}{exo:g8:voltage-voltmeter:4}
\omterm{def:g8:intensity-ammeter:ammeter}{Amperímetro}: \omterm{def:g4:series-circuits:series}{en serie}, con la carretera rota para él, sin
oponerse casi a nada. \omterm{def:g8:voltage-voltmeter:voltmeter}{Voltímetro}: \omterm{def:g5:series-parallel-circuits:parallel}{en paralelo}, sin
romper nada, oponiéndose casi a todo. El error caro: un
\omterm{def:g8:intensity-ammeter:ammeter}{amperímetro} conectado cruzado --- \omterm{def:g5:series-parallel-circuits:parallel}{en paralelo} --- se
convierte en un \omterm{def:g7:short-circuits-safety:device}{cortocircuito} con esfera.
\end{solution}

\begin{solution}{exo:g8:voltage-voltmeter:5}
El empujón para el que está construido un aparato. Con \qty{1.5}{V}: un
resplandor soñoliento. Con \qty{6}{V}: el brillo y la vida de diseño.
Con \qty{12}{V}: un instante brillante y el filamento se parte.
\end{solution}

\begin{solution}{exo:g8:voltage-voltmeter:6}
\omterm{def:g3:first-electric-circuit:battery}{Pila}: \qty{4.5}{V}; lámpara: \qty{4.5}{V}; cable: unos
\qty{0}{V}. El empujón vive en los trabajadores --- el cable es una
carretera llana.
\end{solution}

\begin{solution}{exo:g8:voltage-voltmeter:7}
Lámpara apagada: \qty{0}{V}; \omterm{ex:g3:first-electric-circuit:switch}{interruptor} abierto: los
\qty{4.5}{V} enteros, esperando en el hueco. Con el clic se
intercambian: el \omterm{ex:g3:first-electric-circuit:switch}{interruptor} cerrado baja a casi cero y la lámpara
encendida se queda con los \qty{4.5}{V}.
\end{solution}

\begin{solution}{exo:g8:voltage-voltmeter:8}
Se niega a llevar corriente: como se opone casi del todo, sorbe por
dentro casi nada, así que el desfile y el brillo de la lámpara siguen
como si nadie los mirara.
\end{solution}

\begin{solution}{exo:g8:voltage-voltmeter:9}
Tres \omterm{def:g3:first-electric-circuit:battery}{pilas} en fila: $1.5 + 1.5 + 1.5 = \qty{4.5}{V}$. Una del revés se
opone a las otras dos: $1.5 + 1.5 - 1.5 = \qty{1.5}{V}$ --- el misterio
de la linterna de hace años, ahora en \omterm{def:g8:voltage-voltmeter:voltage}{voltios}.
\end{solution}

\begin{solution}{exo:g8:voltage-voltmeter:10}
Batería recién cargada: un pelín por encima de la nominal ---
brillo pleno y alegre. A medida que la velada vacía la
batería hacia los \qty{11.2}{V}, la lámpara se apaga suavemente
por debajo de su resplandor de diseño. No se quemó nada: la
\omterm{def:g8:voltage-voltmeter:voltage}{tensión} solo flaqueó \emph{por debajo} de la nominal; el contrato
castiga el exceso, no la escasez.
\end{solution}

\begin{solution}{exo:g8:voltage-voltmeter:11}
El \omterm{def:g8:voltage-voltmeter:voltmeter}{voltímetro} plantado en la carretera se opone casi a todo: el
desfile se para prácticamente y la lámpara se queda apagada. El
aparato, que soporta él solo todo el empujón del lazo, lee casi la
\omterm{def:g8:voltage-voltmeter:voltage}{tensión} entera de la \omterm{def:g3:first-electric-circuit:battery}{pila} --- un número correcto para
un circuito mal construido: instructivo, y sin incendio.
\end{solution}

\begin{solution}{exo:g8:voltage-voltmeter:12}
\omterm{def:g5:series-parallel-circuits:parallel}{En paralelo}, cada lámpara abarca directamente la
\omterm{def:g3:first-electric-circuit:battery}{pila}: \qty{4.5}{V} cada una --- empujón nominal completo, brillo
completo, como prometía la vieja regla del paralelo. \omterm{def:g4:series-circuits:series}{En serie}, las
dos lámparas \emph{se reparten} los \qty{4.5}{V} --- unos
\qty{2.25}{V} cada una, una pareja mal alimentada que luce poco, como
observaba la vieja regla de la serie. El capítulo de dentro de dos
convierte ese reparto en una ley de tensiones.
\end{solution}

\begin{solution}{pb:g8:voltage-voltmeter:1}
\textbf{1.} Pinza las dos puntas del \omterm{def:g8:voltage-voltmeter:voltmeter}{voltímetro} a los dos
\omterm{def:g3:first-electric-circuit:battery}{bornes} de la \omterm{def:g3:first-electric-circuit:battery}{pila} --- cruzado, \omterm{def:g5:series-parallel-circuits:parallel}{en paralelo}; un
\omterm{def:g8:voltage-voltmeter:voltmeter}{voltímetro} no rompe nada y no necesita ningún circuito:
mide el empujón que espera.
\textbf{2.} \qty{1.58}{V}: nueva. \qty{1.32}{V}: utilizable.
\qty{0.9}{V}: cansada. \qty{0.1}{V}: muerta --- del cajón a la
basura.
\textbf{3.} Unos $8.9 \div 6 \approx \qty{1.48}{V}$ cada una --- seis
\omterm{def:g3:first-electric-circuit:battery}{pilas} sanas.
\textbf{4.} Su química sigue levantando un empujón debilucho, pero ya
no puede mover un desfile útil: la \omterm{def:g8:voltage-voltmeter:voltage}{tensión} sobrevive como una
promesa que la \omterm{def:g3:first-electric-circuit:battery}{pila} no tiene \omterm{def:g2:pushes-and-pulls:force}{fuerzas} para cumplir bajo carga.
\textbf{5.} Dos en fila: espera $2 \times 1.5 = \qty{3}{V}$; las
supervivientes entregan $2 \times 1.32 = \qty{2.64}{V}$ --- la cámara
puede funcionar, refunfuñando.
\textbf{6.} Las cuatro \omterm{def:g3:first-electric-circuit:battery}{pilas} tienen que sumar \qty{5}{V} o más: cuatro
nuevas de $1.58$ dan \qty{6.3}{V} --- perfecto; mezclando las de
$1.32$, sirve cualquier tripulación que sume al menos \qty{5}{V} (por
ejemplo dos nuevas y dos utilizables: $1.58 \times 2 + 1.32 \times 2 =
\qty{5.8}{V}$). Las cansadas y las muertas que no se presenten.
\textbf{7.} Tres \omterm{def:g3:first-electric-circuit:battery}{pilas}: \qty{4.5}{V} sobre una \omterm{def:g3:first-electric-circuit:bulb}{bombilla} de
\qty{3}{V} --- la mitad otra vez por encima de la nominal: en el mejor
de los casos una velada gloriosa de exceso de brillo, y lo más probable
una quemadura inmediata. El contrato castiga el exceso.
\textbf{8.} Asesinato lento, suave pero real: \qty{9}{V} en un aparato
de \qty{8}{V} es un exceso constante de empujón --- funcionará, se
calentará y envejecerá deprisa. Nominal significa nominal; ``casi'' es
la manera en que mueren jóvenes los filamentos.
\textbf{9.} El empujón de la \omterm{def:g3:first-electric-circuit:battery}{pila} espera, entero, en el único
hueco verdadero --- el \omterm{ex:g3:first-electric-circuit:switch}{interruptor} abierto; la \omterm{def:g3:first-electric-circuit:bulb}{bombilla}
parada, que no lleva corriente, no enseña nada. Después del clic:
\omterm{ex:g3:first-electric-circuit:switch}{interruptor} casi a \qty{0}{V}, \omterm{def:g3:first-electric-circuit:bulb}{bombilla} casi a
\qty{3.1}{V} --- el empujón se traslada del hueco al trabajador.
\textbf{10.} Que el buen cable es una carretera llana: entrega el
desfile sin gastar empujón --- toda la \omterm{def:g8:voltage-voltmeter:voltage}{tensión} se guarda para los
trabajadores.
\textbf{11.} ``El enchufe son \qty{230}{V} --- cincuenta veces la \omterm{def:g3:first-electric-circuit:battery}{pila}
más fuerte de este cajón; cincuenta veces el empujón mueve cincuenta
veces la corriente por las mismas manos húmedas, y el corazón se para
con centésimas de \omterm{def:g8:intensity-ammeter:intensity}{amperio}. Las auditorías terminan en el
cajón''.
\textbf{12.} Por ejemplo: ``Un \omterm{def:g8:voltage-voltmeter:voltmeter}{voltímetro} se pone cruzado, nunca
en la carretera. La \omterm{def:g8:voltage-voltmeter:nominal}{tensión nominal} de una etiqueta promete el
rendimiento de diseño y amenaza con la quemadura si se pasa uno. Y en
cualquier circuito parado, la \omterm{def:g8:voltage-voltmeter:voltage}{tensión} espera en el hueco --- en el
\omterm{ex:g3:first-electric-circuit:switch}{interruptor} abierto, lista''.
\end{solution}
